% Blueprint: Fundamental group
% Created by Numina

\begin{document}

\section{The Fundamental Theorem of Finitely Generated Abelian Groups}

\begin{definition}[Group]
    \label{def:group}
    A \emph{group} is a set $G$ equipped with a binary operation
    $\cdot : G \times G \to G$ satisfying:
    \begin{enumerate}
        \item (Associativity) $(a \cdot b) \cdot c = a \cdot (b \cdot c)$ for all $a, b, c \in G$.
        \item (Identity) There exists $e \in G$ with $e \cdot a = a \cdot e = a$ for all $a \in G$.
        \item (Inverses) For every $a \in G$ there exists $a^{-1} \in G$ with
        $a \cdot a^{-1} = a^{-1} \cdot a = e$.
    \end{enumerate}
\end{definition}

\begin{definition}[Abelian group]
    \label{def:abelian-group}
    \uses{def:group}
    A group $(G, \cdot)$ is \emph{abelian} if $a \cdot b = b \cdot a$ for all $a, b \in G$.
    We typically write the operation additively as $a + b$, with identity $0$ and inverse $-a$.
\end{definition}

\begin{definition}[Subgroup and homomorphism]
    \label{def:subgroup-homomorphism}
    \uses{def:abelian-group}
    A subset $H \subseteq G$ is a \emph{subgroup} if it is closed under the operation and
    inverses and contains $0$. A \emph{homomorphism} $\varphi : G \to H$ of abelian groups is
    a map satisfying $\varphi(a + b) = \varphi(a) + \varphi(b)$. An \emph{isomorphism} is a
    bijective homomorphism, written $G \cong H$.
\end{definition}

\begin{definition}[Direct sum]
    \label{def:direct-sum}
    \uses{def:abelian-group}
    Given abelian groups $A_1, \dots, A_n$, their \emph{direct sum} is
    \[
        A_1 \oplus \cdots \oplus A_n
        = \{ (a_1, \dots, a_n) : a_i \in A_i \},
    \]
    with componentwise addition $(a_i) + (b_i) = (a_i + b_i)$.
\end{definition}

\begin{definition}[Cyclic group]
    \label{def:cyclic-group}
    \uses{def:abelian-group, def:subgroup-homomorphism}
    An abelian group $G$ is \emph{cyclic} if there exists $g \in G$ such that
    $G = \{ ng : n \in \mathbb{Z} \}$. Every cyclic group is isomorphic to either $\mathbb{Z}$
    (infinite order) or $\mathbb{Z}/n\mathbb{Z}$ for some $n \ge 1$ (finite order $n$).
\end{definition}

\begin{definition}[Finitely generated]
    \label{def:finitely-generated}
    \uses{def:abelian-group}
    An abelian group $G$ is \emph{finitely generated} if there exist $g_1, \dots, g_n \in G$
    such that every element can be written as $k_1 g_1 + \cdots + k_n g_n$ with
    $k_i \in \mathbb{Z}$.
\end{definition}

\begin{definition}[Torsion and free part]
    \label{def:torsion-free-part}
    \uses{def:abelian-group, def:subgroup-homomorphism}
    An element $g \in G$ is a \emph{torsion element} if $ng = 0$ for some $n \ge 1$. The set
    of torsion elements forms a subgroup $T(G) \subseteq G$, called the \emph{torsion subgroup}.
    A group is \emph{torsion-free} if $T(G) = 0$, and \emph{free abelian of rank $r$} if it is
    isomorphic to $\mathbb{Z}^r$.
\end{definition}

\begin{theorem}[Fundamental Theorem, Invariant Factor Form]
    \label{thm:invariant-factor-form}
    \uses{def:abelian-group, def:direct-sum, def:cyclic-group, def:finitely-generated,
        def:torsion-free-part, def:subgroup-homomorphism}
    Let $G$ be a finitely generated abelian group. Then there exist a unique integer
    $r \ge 0$ and unique integers $1 < d_1 \mid d_2 \mid \cdots \mid d_k$ such that
    \[
        G \;\cong\; \mathbb{Z}^r \;\oplus\; \mathbb{Z}/d_1\mathbb{Z} \;\oplus\;
        \mathbb{Z}/d_2\mathbb{Z} \;\oplus\; \cdots \;\oplus\; \mathbb{Z}/d_k\mathbb{Z}.
    \]
    The integer $r$ is the \emph{rank} (or \emph{Betti number}) of $G$, and
    $d_1, \dots, d_k$ are the \emph{invariant factors}.
\end{theorem}
\begin{proof}
    \uses{def:torsion-free-part, def:direct-sum, def:finitely-generated, def:cyclic-group}
    The torsion subgroup $T(G)$ of a finitely generated abelian group is itself finitely
    generated, hence finite, and the quotient $G / T(G)$ is finitely generated and torsion-free,
    hence free abelian of some rank $r \ge 0$. The short exact sequence $0 \to T(G) \to G
    \to G/T(G) \to 0$ splits because $G/T(G)$ is free, giving $G \cong \mathbb{Z}^r \oplus T(G)$.
    Applying the structure theorem for finitely generated modules over the principal ideal
    domain $\mathbb{Z}$ (via Smith normal form) to $T(G)$ produces a decomposition
    $T(G) \cong \mathbb{Z}/d_1\mathbb{Z} \oplus \cdots \oplus \mathbb{Z}/d_k\mathbb{Z}$ with
    $1 < d_1 \mid d_2 \mid \cdots \mid d_k$. Uniqueness of $r$ and of the invariant factors
    follows from the uniqueness part of Smith normal form.
\end{proof}

\begin{theorem}[Fundamental Theorem, Elementary Divisor Form]
    \label{thm:elementary-divisor-form}
    \uses{def:abelian-group, def:direct-sum, def:cyclic-group, def:finitely-generated,
        def:torsion-free-part, def:subgroup-homomorphism}
    Let $G$ be a finitely generated abelian group. Then there exist an integer $r \ge 0$,
    primes $p_1, \dots, p_s$ (not necessarily distinct), and positive integers
    $e_1, \dots, e_s$ such that
    \[
        G \;\cong\; \mathbb{Z}^r \;\oplus\; \bigoplus_{i=1}^{s} \mathbb{Z}/p_i^{e_i}\mathbb{Z},
    \]
    and the multiset of prime powers $\{p_i^{e_i}\}$ is uniquely determined by $G$. These are
    the \emph{elementary divisors} of $G$.
\end{theorem}
\begin{proof}
    \uses{thm:invariant-factor-form, lem:passage-between-forms}
    Start from the invariant factor decomposition
    $G \cong \mathbb{Z}^r \oplus \mathbb{Z}/d_1\mathbb{Z} \oplus \cdots
    \oplus \mathbb{Z}/d_k\mathbb{Z}$ given by \cref{thm:invariant-factor-form}. By
    \cref{lem:passage-between-forms}, each factor $\mathbb{Z}/d_j\mathbb{Z}$ splits as a
    direct sum $\bigoplus_\ell \mathbb{Z}/p_{j,\ell}^{a_{j,\ell}}\mathbb{Z}$ over the prime
    factorization of $d_j$. Concatenating these decompositions gives the claimed elementary
    divisor form. Uniqueness of the multiset of prime powers follows from the uniqueness of
    invariant factors together with the uniqueness of the prime factorization of each $d_j$.
\end{proof}

\begin{lemma}[Passage between the two forms]
    \label{lem:passage-between-forms}
    \uses{def:direct-sum, def:cyclic-group}
    If $d = p_1^{a_1} \cdots p_m^{a_m}$ is a prime factorization with the $p_i$ distinct
    primes and each $a_i \ge 1$, then
    \[
        \mathbb{Z}/d\mathbb{Z} \;\cong\; \mathbb{Z}/p_1^{a_1}\mathbb{Z} \;\oplus\; \cdots
        \;\oplus\; \mathbb{Z}/p_m^{a_m}\mathbb{Z}.
    \]
    Splitting each invariant factor in this way yields the elementary divisors; conversely,
    grouping the elementary divisors by taking, for each prime, the largest power into $d_k$,
    the next largest into $d_{k-1}$, and so on, recovers the invariant factors.
\end{lemma}
\begin{proof}
    \uses{def:direct-sum, def:cyclic-group}
    The isomorphism is a direct application of the Chinese Remainder Theorem: since the
    prime powers $p_i^{a_i}$ are pairwise coprime and their product equals $d$, the natural
    ring homomorphism
    $\mathbb{Z}/d\mathbb{Z} \to \prod_{i=1}^m \mathbb{Z}/p_i^{a_i}\mathbb{Z}$
    is an isomorphism of rings, and in particular of abelian groups. Reversing the grouping
    as described recovers the original invariant factors because, for each prime $p$, the
    divisibility chain $d_1 \mid d_2 \mid \cdots \mid d_k$ forces the $p$-parts of the $d_j$
    to form an increasing sequence of prime powers.
\end{proof}

\end{document}
